\subsubsection{Lucas theorem}
\begin{lstlisting}[style=mystyle, language=C++]
// TC: O(log_p n) per query
// usa: C(n, k) mod p con p primo y n, k grandes
#include <bits/stdc++.h>
using namespace std;

// usa: C(n, k) mod p cuando n y k son gigantes y p es primo
// idea: escribir n y k en base p; C(n, k) congruente Productorio C(ni, ki) mod p
//       (un producto por cada digito en base p)
// requires fact[] and invFact[] already built for MOD (small case)
long long nCrLucas(long long n, long long k, long long MOD){
	if(k == 0) return 1;
	long long ni = n % MOD, ki = k % MOD;
	if(ki > ni) return 0;
	return nCrLucas(n / MOD, k / MOD, MOD) * nCr((int)ni, (int)ki) % MOD;
}

int main(){
	buildFact(MOD);  // precompute factorials mod MOD
	// C(1000, 100) mod 998244353 using Lucas
	cout << nCrLucas(1000, 100, MOD) << "\n";
	return 0;
}
\end{lstlisting}
